\chapter{Framework}\label{ch:setup}

\paragraph{Tasks and demonstrations.}
We use the standard noisy in-context linear regression setting \citep{garg2022,raventos2023}. A task is a vector $w\in\R^d$ drawn from the baseline prior $\mathcal{N}(0,a_0 I_d)$. A prompt contains $n$ demonstrations $D_n=\{(x_i,y_i)\}_{i=1}^n$ and a query $x$, where
\[
x_i,x\overset{\mathrm{iid}}{\sim}\mathcal{N}(0,I_d),\qquad
y_i=w^\top x_i+\epsilon_i,\quad
\epsilon_i\overset{\mathrm{iid}}{\sim}\mathcal{N}(0,1).
\]
The parameter $a_0>0$ controls variation across tasks. Numerical illustrations in Chapters~\ref{ch:specificity}--\ref{ch:reliability} use $d=16$ and $a_0=10$.

\paragraph{Descriptions as uncertain task information.}
A description supplies a center $m\in\R^d$. Its \emph{specificity} is a precision $c\geq 1/a_0$, and its \emph{reliability} is a probability $p\in[0,1]$. Write $r=(ca_0)^{-1}\in(0,1]$ for the fraction of baseline variance remaining when the description is relevant. With an unobserved relevance indicator $Z\sim\mathrm{Bernoulli}(p)$, independent of $m$, we generate
\begin{align*}
m&\sim\mathcal{N}(0,(1-r)a_0I_d),\\
w\mid m,Z=1&\sim\mathcal{N}(m,ra_0I_d),\\
w\mid m,Z=0&\sim\mathcal{N}(0,a_0I_d).
\end{align*}
Thus the description induces the mixture prior
\[
w\mid m\sim p\,\mathcal{N}(m,c^{-1}I_d)
 +(1-p)\,\mathcal{N}(0,a_0I_d),
\]
as in robust mixture priors for potentially incompatible information \citep{schmidli2014}. Both relevance branches have the same marginal task distribution, $w\sim\mathcal{N}(0,a_0I_d)$. Consequently, the comparison with and without a description does not change the task population.

A description is \emph{reliable} when $p=1$ and \emph{unreliable} when $p<1$. Greater $c$ means greater specificity conditional on relevance. The model separates uncertainty within the described region from uncertainty about whether that region applies. It does not assume that every natural-language description can be summarized by these two scalars.

A relevant description splits the baseline variance $a_0$ into the part $(1-r)a_0$ that $m$ resolves and the part $ra_0$ that it leaves. Because the observation noise has unit variance, $c$ is measured in demonstrations: in each direction, it is the ESS that \citet{morita2008} assign to a normal prior. We call a description \emph{loose} when $c$ is small relative to $d$ and \emph{specific} when $c$ exceeds $d$. At $d=16$ and $a_0=10$, a description with $c=1$ ($r=0.1$) is loose, and one with $c=100$ ($r=0.001$) is specific.

\paragraph{Prediction error.}
For a prediction $\hat y$ of the query response, we measure error against its conditional mean:
\begin{equation}\label{eq:regret}
R=\E\big[(\hat y-w^\top x)^2\big].
\end{equation}
This excludes irreducible query noise. The expectation is over tasks, descriptions, demonstrations, and queries. For a fixed learner, let $R^{\mathrm{desc}}(n)$ and $R^{\mathrm{demo}}(n)$ denote its errors with and without a description. We write $R^{\mathrm{desc}}_{p=1}(n)$ for the reliable case.

\begin{definition}[Predictive effective sample size]\label{def:ess}
Given $n$ demonstrations, let $n^*$ be the smallest $t\geq0$ such that $R^{\mathrm{demo}}(t)=R^{\mathrm{desc}}(n)$. When this crossing exists, the ESS of the description is
\[
\ess_n=n^*-n.
\]
We interpolate $R^{\mathrm{demo}}$ between available demonstration counts: linearly in $t$ for the theory and meta-trained models, and in $\log t$ for language models. For measured curves, we use the first crossing within the observed range and report out-of-range values as bounds. The language-model conventions are detailed in Section~\ref{app:llm}. If $R^{\mathrm{desc}}(n)>R^{\mathrm{demo}}(t)$ for all $t\geq0$, we report $\ess_n\leq-n$.
\end{definition}

ESS measures a horizontal displacement of a learning curve: $\ess_n=10$ means that $n$ demonstrations plus the description match the error of $n+10$ demonstrations without it. For a decreasing baseline curve, a negative value indicates increased error at the same demonstration count. Each model is evaluated against its own curve, so ESS measures the description's value \emph{to that learner}; a large ESS alone does not imply low absolute error. Under the correctly specified Bayesian model, descriptions cannot increase expected risk, and ESS is nonnegative.

Because $R$ averages over descriptions, the ESS belongs to a kind of description, determined by $c$ and $p$, rather than to a particular value of $m$.
